\begin{table}[H]
\centering
\caption{Productividad y remuneración por departamento, 2024pr}
\label{tab:dept_productividad_remuneracion_33}
\scriptsize
\begin{tabular}{lrrrrr}
\toprule
Departamento & PIB/hora 2024pr & Rem./hora 2024 & Residuo & PIB/trab. 2024pr & Rem./trab. 2024 \\
\midrule
Bogotá D.C. & 28,6 & 15,0 & 79,1\% & 67,4 & 2,95 \\
Vaupés & 26,2 & 12,1 & 45,3\% & 58,5 & 2,26 \\
Antioquia & 20,3 & 10,9 & 31,9\% & 47,7 & 2,12 \\
Guainía & 17,0 & 9,7 & 19,2\% & 38,8 & 1,85 \\
Caldas & 15,6 & 9,4 & 15,9\% & 36,6 & 1,85 \\
Casanare & 67,8 & 10,5 & 15,7\% & 165,2 & 2,13 \\
Risaralda & 16,4 & 9,2 & 13,2\% & 38,7 & 1,82 \\
San Andrés y Providencia & 33,3 & 9,6 & 13,1\% & 82,4 & 1,97 \\
Cundinamarca & 17,4 & 9,2 & 12,8\% & 40,4 & 1,79 \\
Quindío & 15,1 & 9,1 & 12,1\% & 33,8 & 1,71 \\
Valle del Cauca & 23,4 & 9,3 & 11,8\% & 54,5 & 1,79 \\
Guaviare & 17,0 & 9,1 & 11,2\% & 39,1 & 1,74 \\
Santander & 26,1 & 9,1 & 9,2\% & 62,0 & 1,80 \\
Meta & 25,9 & 9,0 & 8,4\% & 63,1 & 1,83 \\
Boyacá & 22,6 & 8,5 & 2,7\% & 47,5 & 1,49 \\
Amazonas & 21,7 & 8,3 & 1,0\% & 50,2 & 1,61 \\
Atlántico & 16,0 & 8,1 & -0,5\% & 37,3 & 1,58 \\
Vichada & 38,1 & 8,0 & -6,7\% & 92,2 & 1,60 \\
Tolima & 16,9 & 7,6 & -7,2\% & 40,0 & 1,50 \\
Caquetá & 10,5 & 7,4 & -7,8\% & 24,9 & 1,47 \\
Putumayo & 202,1 & 10,4 & -8,4\% & 443,0 & 1,90 \\
Norte de Santander & 10,0 & 7,4 & -8,5\% & 24,3 & 1,49 \\
Huila & 14,2 & 7,3 & -9,7\% & 31,6 & 1,36 \\
Bolívar & 19,8 & 7,0 & -14,9\% & 45,1 & 1,33 \\
Nariño & 9,1 & 6,5 & -19,2\% & 17,8 & 1,06 \\
Cauca & 13,7 & 6,5 & -19,5\% & 26,6 & 1,06 \\
Cesar & 13,8 & 6,4 & -20,8\% & 32,6 & 1,26 \\
Córdoba & 10,7 & 6,1 & -24,7\% & 22,0 & 1,04 \\
Magdalena & 10,3 & 5,9 & -26,2\% & 23,9 & 1,15 \\
Chocó & 11,9 & 5,8 & -28,2\% & 28,6 & 1,16 \\
Sucre & 10,4 & 5,7 & -29,2\% & 22,0 & 1,01 \\
Arauca & 66,8 & 6,1 & -31,9\% & 171,2 & 1,31 \\
La Guajira & 11,1 & 4,8 & -40,6\% & 24,8 & 0,89 \\
\bottomrule
\end{tabular}
\caption*{\footnotesize Nota: PIB por hora en miles de pesos constantes de 2015; remuneración por hora en miles de pesos constantes de 2025; PIB por trabajador en millones de pesos constantes de 2015; remuneración por trabajador en millones de pesos constantes de 2025 al mes. El residuo mide la distancia porcentual frente a la tendencia lineal entre PIB por hora y remuneración por hora. Fuente: cálculos propios con DANE y GEIH.}
\end{table}